\begin{graphicspathcontext}{{./chapters/ai/imgs/},{./chapters/ai/imgs/auto/},\old}

\adaptAIdefinitioncontent
\talkadaptivesidecite
\begin{frame}[c]{Qu'est-ce que l'IA ?}
	\talkadaptivesize
	\ifMinskyAI{
		\begin{definition}[Minsky 1961]
			l'IA est la science de \Emph{faire faire aux machines des choses} qui \emph{exigeraient de l'intelligence si elles étaient faites par des humains}
		\end{definition}
	}\fi
	\talkadaptivevspacea
	\ifOECDAI{
		\begin{definition}[OCDE 2024]
			Une IA est un système fondé sur des machines qui, pour des \Emph{objectifs explicites ou implicites}, déduit à partir de \Emph{données} qu'il reçoit comment générer des résultats tels que des \Emph{prévisions, du contenu, des recommandations ou des décisions} susceptibles d'influencer les environnements physiques ou virtuels. Les \Emph{différents systèmes d'IA} varient dans leurs niveaux d'autonomie et d'adaptabilité une fois déployés
		\end{definition}
	}\fi
	\talkadaptivevspaceb
	\ifEUAIAct{
		\begin{definition}[EU IA Act 2024]
			Un \Emph{système d'IA} est un système à base machine qui est conçu pour fonctionner avec différents niveaux d'\Emph{autonomie} et qui peut faire preuve d'\Emph{adaptabilité} après son déploiement, et qui, pour des \Emph{objectifs explicites ou implicites}, déduit, à partir de \Emph{données} qu'il reçoit, la manière de générer des résultats tels que des \Emph{prédictions, du contenu, des recommandations ou des décisions} pouvant influencer les environnements physiques ou virtuels
		\end{definition}
	}\fi
\end{frame}

\end{graphicspathcontext}

\endinput

